\clearpage\thispagestyle{empty}
\pdfbookmark[0]{Quellen zum Ceres-Bericht vom Oktober 1802}{ceres180210-sources}
\begin{center}\textbf{Quellen zum Ceres-Bericht vom Oktober 1802}\end{center}
\noindent Diese KI-Übertragung in heutiges Deutsch folgt dem Auszug aus Franz Xaver von Zachs Bericht über Ceres in \href{https://gdz.sub.uni-goettingen.de/id/PPN236007467}{Gauß’ \textit{Werken}, Band~VI} (Göttingen, 1874), S.~224--225. Von Zach zitiert Johann Friedrich Wurm und stellt Carl Friedrich Gauß’ Berechnung der Störungen der Ceres durch Jupiter vor. Mit dem lateinischen \textit{quid pro quo} meint Wurm eine Verwechslung.
\par\smallskip\noindent Der \href{https://api.digitale-sammlungen.de/iiif/image/v2/bsb10538598_00392/full/1450,/0/default.jpg}{Erstdruck vom Oktober 1802} steht als Beitrag XXXVIII in der \textit{Monatlichen Correspondenz}, Band~VI, S.~382--389. Der Auszug in den \textit{Werken} entspricht dem \href{https://api.digitale-sammlungen.de/iiif/image/v2/bsb10538598_00397/full/1450,/0/default.jpg}{Text und den Jahreswerten auf S.~387} und den \href{https://api.digitale-sammlungen.de/iiif/image/v2/bsb10538598_00398/full/1450,/0/default.jpg}{periodischen Gleichungen auf S.~388}.
\par\smallskip\noindent Der Erstdruck bezeichnet \(0{,}000005909\) als jährliche Zunahme der Exzentrizität. In den \textit{Werken} heißt es jährliche Abnahme mit dem Wert \(0.000005909\). Das Argument des Längenterms \(+56{,}69\) endet in der Zeitschrift mit \(11^\circ\ 24'\ 1'\), die entsprechende Angabe in den \textit{Werken} lautet \(11\ 24\ 1''\). Beim Term \(+0{,}0002469\) des Radiusvektors steht in der Zeitschrift \(23^\circ\ 42'\ 43''\), in den \textit{Werken} \(23\ 42\ 53''\). Beide Übertragungen folgen den \textit{Werken}.
\par\smallskip\noindent Übertragung und Satz wurden mit KI erstellt. Die Digitalisate der Zeitschrift stammen von der Bayerischen Staatsbibliothek.
\clearpage
